% !TeX root = ../抽象代数笔记.tex
% 由主文档通过 \input 载入；请编译主文档。
\setcounter{chapter}{10} % 保留原书章号。
\chapter{理想、商环与中国剩余定理}\label{ch:ideals}

\begin{description}
  \item[学习目标] 理解理想控制商环运算，能够用极大理想构造域并分解适当的商环。
  \item[先修知识] 第\ref{ch:rings}章及群同构定理。
\end{description}

\section{理想与生成理想}
定义左理想、右理想和双边理想（ideal）；构造商环时均使用双边理想。说明理想是加法子群并满足乘法吸收条件，不能与子环混为一谈。对交换环定义主理想 $(a)$、有限生成理想 $(a_1,\ldots,a_r)$，介绍理想的和、交与积。

\section{商环及环同构定理}
对环 $R$ 的双边理想 $I$，验证 $R/I$ 的加法与乘法良定义。证明环同态的核为双边理想，以及 $R/\ker f\cong\operatorname{im}f$。给出对应定理和迭代取商公式，并用求值映射展示“关系式被设为零”的商环解释。

\section{素理想与极大理想}
本节设 $R$ 为交换环，且所讨论理想均为真理想。定义素理想（prime ideal）与极大理想（maximal ideal），证明 $P$ 为素理想当且仅当 $R/P$ 为整环，$M$ 为极大理想当且仅当 $R/M$ 为域。推出极大理想必为素理想，并用 $\mathbb Z$ 的零理想说明逆命题不成立。

\section{中国剩余定理}
对交换环 $R$ 中两两互素的有限个理想 $I_1,\ldots,I_r$，其中互素指 $I_i+I_j=R$（$i\ne j$），证明自然同态满射且其核为各理想之交，并证明此时交等于积，从而得到
\begin{equation}\label{eq:ring-crt}
  R/(I_1\cdots I_r)\cong\prod_{j=1}^{r}R/I_j.
\end{equation}
先处理两个理想，再归纳推广；对整数模互素数的情形，给出由贝祖等式构造原像的方法。

\section{例题、练习与自检}
解一组模 $3$、$5$、$7$ 的联立同余，分析 $\mathbb Z/60\mathbb Z$ 的直积分解。比较 $\mathbb R[x]/(x^2+1)$ 与 $\mathbb R[x]/(x^2-1)$，说明商环分别表现为域与含零因子的环。练习要求写出同构的具体映射，并指出理想不互素时式\eqref{eq:ring-crt}的满射性为何可能失败。

